\Def Декартовым деревом называют бинарное дерево, содержащее в
себе пары $(x_i,\,y_i)$, при этом данное дерево является деревом поиска по
ключам (то есть по значениям $x_i$) и кучей по приоритетам (то есть по
значениям $y_i$).

\Th Для любого набора пар ключ – приоритет декартово
дерево существует. Более того, если все приоритеты различны, то
декартово дерево единственно.

\Proof Рассмотрим графическую интерпретацию. Выбираем в
качестве корня $(x_0, y_0)$ узел с максимальным приоритетом, если
таковых несколько, то с медианным ключом среди таковых.Получаем
два набора, образующих левое поддерево: все те, у кого ключи $\leq\,x_0$, и
правое: все те, у кого ключи > $x_0$.Таким образом, рекурсивно перешли
к наборам меньше, при этом для дерева из одного элемента все
корректно.
Теперь единственность. Обратимся к алгоритму выше и заметим, что
на каждом шаге корень поддерева определен однозначно. \Endproof

\Thbd В декартовом дереве из N узлов, приоритеты y которого являются случайными величинами c равномерным распределением, средняя глубина вершины $O(log\,N)$.

\subsection*{Операции Split, Merge}

Рассмотрим операцию Merge. Напомним, что все ключи $T_1$ меньше
ключей $T_2$.

\begin{itemize}
    \item Если $T_1.root.y\,>\,T_2.root.y\text{, то }T.root\,=\,T_1.root$.
    \item Значит известно левое поддерево T: $T.root.left\,=\,T_1.root.left$.
    \item $T.root.right\,=\,Merge(T_1.root.right, T_2)$
    \drawsomemedium{pix/treap_merge.png}
    \item Иначе симметрично действуем.
\end{itemize}

\pagebreak
Рассмотрим операцию Split

\begin{itemize}
    \item Если $x > T.root.x$, то известно левое поддерево $T_1$: $T_1.left\,=\,T.left$.
    \item Тогда $\langle T_1.right,\,T_2\rangle\,=\,Split(T.right,\,x)$  
    \drawsomemedium{pix/treap_split.png}
    \item Иначе симметрично действуем.
\end{itemize}

После операции Split будут $T_1$ и $T_2$ такие, что в $T_1$ находятся все ключи меньшие $x$, а в $T_2$ больше, либо равные.

\subsection*{Выражение операций Insert и Erase через Merge и Split}

Рассмотрим операцию Insert. Дано декартово дерево T ключ K. Нужно вставить ключ в дерево

\begin{enumerate}
    \item Если K уже есть в T, то алгоритм окончен.
    \item Иначе делаем $Split(T,\,x)\,\rightarrow \, \langle T_1,\,T_2\rangle$.
    \item Возвращаем результат $Merge(Merge(T_1,\,Tree(x)),\,T_2)$.
\end{enumerate}

Рассмотрим операцию Erase. Дано декартово дерево T и ключ K. Нужно удалить узел с ключом K, если такой есть.

\begin{enumerate}
    \item Сначала спускаемся по дереву (как в обычном бинарном дереве поиска), ища удаляемый элемент
    \item Найдя элемент, вызываем слияние его левого и правого сыновей
    \item Результат Merge ставим на место удаляемого элемента
\end{enumerate}

\Statement Время работы Erase, Insert составляет $O(log\,N)$ в среднем

\pagebreak